\chapter{おわりに}

この本を手に取って下さり、ありがとうございました。
今回は初のサークル参加で、時間もなく、少し説明が雑な部分が多くなってしまった気もしますが、論理回路に少しでも興味をもらえれば幸いです。
本書では論理回路の仕組みを中心に話をしてきましたが、割愛した部分も多々あります。今後サークルとしてまた本を出す際には今回端折ってしまった所も説明できたらいいと思っています。
今後も電子回路やロボコン関する本を出していく予定なので、応援してもらえると嬉しいです。

by \quad Ari3
\\
\\

私はもともとコミケなど名前しか知らない状態で、全く関係ない人生を送っていたのですが、Ari3に無理矢理巻き込まれていつの間にか本を書くことになってしまいました。
大学では電気電子工学を専攻しているため、論理回路の理解自体はそこまで難しくなかったのですが、文章にするのは非常に難しく、特に図の作成には苦労しました。
また、Ari3に本を書けと言われたときはすでに12月で、非常にタイトなスケジュールであったため、書きたいこともすべて書ききれず、説明も雑になってしまった部分も多々あるかと思います。
しかし、論理回路の基礎を理解する上で最低限必要なことは盛り込めたと思うので、論理回路に興味を持ってもらえれば幸いです。
来年はより深堀した内容の本を書けたらいいなと思っています。

by \quad El Presidente


\section*{参考文献}
\begin{itemize}
    \item 『ディジタル回路』、斉藤 忠夫、株式会社コロナ社、昭和57年
    \item 『ディジタル回路』、五島 正裕、株式会社数理工学社、2007年
    \item 『ディジタル電子回路』、木村 誠聡、株式会社数理工学社、2012年
\end{itemize}

\newpage
\thispagestyle{empty}

\vspace*{\fill}

\vspace{15em}

\noindent\rule{\textwidth}{0.4pt}

\vspace{1em}

\begin{center}
{\bfseries 理解れ! 論理回路入門}
\end{center}

\vspace{1em}

\noindent
\begin{tabular}{@{}ll@{}}

著者 & Ari3 \quad \& \quad El Presidente \\
サークル & ありのす工房 \\ 
発行日 & 2025年12月31日 初版\\
印刷所 & ペンタロー\\
組版 & \LaTeX \\
連絡先 & chaume.yu@gmail.com\\
twitter &  @early72412338
\end{tabular}

\vspace{1em}

\noindent\rule{\textwidth}{0.4pt}

